\documentclass[11pt]{article}
\usepackage{amsmath,amssymb,amsthm,geometry}
\geometry{margin=1in}
\newtheorem{theorem}{Theorem}[section]
\newtheorem{lemma}[theorem]{Lemma}
\begin{document}
% No preamble. Unique prefix oi-. Standard theorem environments and amsmath.
% Abstract overlapping-support Gaussian cutoff theorem; actual RG matching
% and strict background support are explicitly separate.

\section{Overlapping interactions in a Gaussian cutoff integral}
\label{oi-section}

The disjoint-cell cutoff expansion extends to overlapping interaction
supports, including supports of unbounded size, under an explicit
weighted multiplier norm. The essential step is to retain the original
cell endpoints when differentiating covariance edges between larger
interaction molecules. This controls the degree factorials without
enlarging the fixed Gaussian cell dimension. The statement below is
abstract; it does not verify the fluctuation coordinates or interaction
bounds of the Yang--Mills construction.

\subsection{Assumptions and quantitative statement}

Let $\Lambda$ be a finite metric set of cells, each carrying
$x_i\in\mathbb R^D$, with $D$ fixed. Let $A$ vary in an open complex
background domain. The complex symmetric covariance $C(A)$ is
holomorphic and obeys, uniformly in volume and background,
\begin{equation}\label{oi-covariance}
 \operatorname{Herm}C\ge\lambda I,\qquad
 \|C\|\le\Lambda_0,\qquad 0<\lambda\le\Lambda_0.
\end{equation}
For $i\ne j$ assume
\begin{equation}\label{oi-decay}
 \|C_{ij}\|\le J_{ij}=J_{ji},\qquad J_{ii}=0,\qquad
 A_\kappa=\sup_i\sum_j e^{\kappa d(i,j)}J_{ij}<\infty.
\end{equation}
Write $\ell(S)$ for the minimum metric spanning-tree length of a
nonempty set $S$, with $\ell(\{i\})=0$. Gaussian expectations use
\[
 p_C(x)=(2\pi)^{-D|\Lambda|/2}(\det C)^{-1/2}
                     \exp(-x^TC^{-1}x/2),
\]
with the analytic determinant branch on the accretive cone.

Fix the real cutoff
\[
 \chi_i=\mathbf1_{\{\|x_i\|_2\le R\}},\qquad
 \chi_Y=\prod_{i\in Y}\chi_i.
\]
Consider a finite labelled family of nonempty interaction supports
$Y\subset\Lambda$. Distinct labels may share a support.
The functions $V_Y(A,x_Y)$ are measurable in the real variables and
holomorphic in $A$ where $\chi_Y=1$, and
\begin{equation}\label{oi-multiplier}
 v_Y=\sup_{A,\;\chi_Y=1}|e^{V_Y(A,x_Y)}-1|<\infty.
\end{equation}
No values outside that cutoff region or real-field derivatives are
required. Define $w_Y=\chi_Y(e^{V_Y}-1)$, extended by zero.
Set
\begin{align}
 \delta&=\lambda/2,&
 V_*&=8\Lambda_0^2/\lambda,&
 K_G&=(2V_*/\lambda)^{D/2},\notag\\
 p_R&=2^{D/2}e^{-R^2/(4V_*)},&
 \tau&=K_Gp_R,& c_G&=e^2/\delta,                 \label{oi-constants}\\
 \epsilon&=\sup_i\sum_{Y\ni i}
      v_Y K_G^{|Y|}e^{4|Y|+\kappa\ell(Y)},&
 \eta&=e^3\tau+\epsilon,&
 \eta_*&=\frac1{4\max(1,c_GA_\kappa)}.           \label{oi-smallness}
\end{align}
The interaction sum counts labels. These sufficient constants are
not asserted to be optimal.

\begin{theorem}[Overlapping Gaussian cutoff interactions]
\label{oi-theorem}
If $\eta\le\eta_*$, then
\[
 Z_V(A)=\mathbb E_{C(A)}
   \left[\prod_{i\in\Lambda}\chi_i
              \exp\left(\sum_YV_Y\right)\right]
\]
is nonzero and has a holomorphic logarithm defined by a convergent
expansion
\begin{equation}\label{oi-log}
 \log Z_V(A)=\sum_{\varnothing\ne X\subset\Lambda}U_V(X;A),\qquad
 \sup_i\sum_{X\ni i}e^{\kappa\ell(X)+|X|}
       \|U_V(X;\cdot)\|_\infty\le2\eta .
\end{equation}
The constants do not depend on the number or sizes of supports.
If
\begin{equation}\label{oi-quarter}
 e^3\tau\le\eta_*/4,\qquad \epsilon\le\eta_*/4,
\end{equation}
the canonical expansion also satisfies
\begin{equation}\label{oi-difference}
 \sup_i\sum_{X\ni i}e^{\kappa\ell(X)+|X|}
      \|U_V(X;\cdot)-U_0(X;\cdot)\|_\infty
 \le8\epsilon ,
\end{equation}
where $U_0$ is the pure Gaussian cutoff forest/Mayer expansion.
The support label $X$ concerns fluctuation cells; strict locality
in the background is a separate assertion.
\end{theorem}

\subsection{First expansion: disjoint interaction molecules}

Put $q_i=-\mathbf1_{\mathrm{bad},i}$. The identity
\begin{equation}\label{oi-product}
 \prod_i\chi_i e^{\sum_YV_Y}
       =\prod_i(1+q_i)\prod_Y(1+w_Y)
\end{equation}
is exact, and $q_iw_Y=0$ whenever $i\in Y$.
For a nonempty set $S$ let
\begin{align}
 a_S(x_S)&=
 \sum_{\substack{\mathcal F\ne\varnothing,\;\cup_{Y\in\mathcal F}Y=S\\
                  \mathcal F\ {\rm overlap\text{-}connected}}}
                         \prod_{Y\in\mathcal F}w_Y(x_Y),\notag\\
 M_S&=
 \sum_{\substack{\mathcal F\ne\varnothing,\;\cup_{Y\in\mathcal F}Y=S\\
                  \mathcal F\ {\rm overlap\text{-}connected}}}
                         \prod_{Y\in\mathcal F}v_Y .       \label{oi-molecules}
\end{align}
Each label is used at most once in these sums. Connectivity refers
to the intersection graph of the chosen supports. Then
$|a_S|\le M_S$. Define $f_S=a_S$ for nonsingletons and
$f_{\{i\}}=a_{\{i\}}+q_i$. Grouping overlap components in
\eqref{oi-product} gives exactly
\begin{equation}\label{oi-disjoint}
 \prod_i\chi_i e^{\sum_YV_Y}
   =\sum_{\mathcal P\ {\rm pairwise\ disjoint}}
                               \prod_{S\in\mathcal P}f_S(x_S).
\end{equation}
The grouping is unique: interaction components and bad singleton
marks have disjoint supports. Introduce
\begin{equation}\label{oi-b}
 b(S)=K_G^{|S|}M_S+\tau\mathbf1_{\{|S|=1\}}.
\end{equation}

\begin{lemma}[Overlap-animal estimate]\label{oi-animal}
If $\epsilon\le1$, then
\begin{equation}\label{oi-primitive-norm}
 \sup_i\sum_{S\ni i}b(S)e^{\kappa\ell(S)+3|S|}
       \le e^3\tau+\epsilon=\eta .
\end{equation}
\end{lemma}
\begin{proof}
Set $u_Y=v_YK_G^{|Y|}e^{3|Y|+\kappa\ell(Y)}$.
For an overlap-connected family,
$\ell(\cup Y)\le\sum_Y\ell(Y)$ and
$|\cup Y|\le\sum_Y|Y|$: join spanning trees of the supports,
which share vertices along their intersection graph.
Since $K_G\ge1$, its contribution to the interaction part of
\eqref{oi-primitive-norm} is at most $\prod_Yu_Y$.

Majorize a connected family anchored at $i$ by rooted spanning
trees of its intersection graph, choosing a root support containing
$i$. Dropping distinctness of labels only increases the sum.
The nonnegative rooted-tree majorant is the monotone iterative
solution of
\[
 T_Y=u_Y\exp\left(\sum_{Z:\,Z\cap Y\ne\varnothing}T_Z\right).
\]
If $T_Z\le u_Ze^{|Z|}$, then
\[
 \sum_{Z:\,Z\cap Y\ne\varnothing}T_Z
 \le\sum_{j\in Y}\sum_{Z\ni j}u_Ze^{|Z|}
 \le\epsilon |Y|\le |Y|.
 \]
Thus the bound persists through iteration, and
$\sum_{Y\ni i}T_Y\le\epsilon$. Every original connected family
is included at least once. The remaining bad singleton term is
$e^3\tau$.
\end{proof}

\subsection{Covariance forests with original cell endpoints}

Fix disjoint primitive supports $S_1,\ldots,S_n$. Introduce a
covariance parameter $s_{ab}$ for each pair of molecules, multiplying
all blocks between $S_a$ and $S_b$ and leaving internal blocks
unchanged. Partition masks in these parameters are direct sums of
actual principal submatrices of $C$. Their convex hull retains
\eqref{oi-covariance}, so the BKAR forest formula applies there.
At a forest mask both the covariance and the product of primitive
functions factorize over the forest components.

For a molecule-tree $T$, expand each covariance derivative into
its original endpoint cells:
\begin{equation}\label{oi-cell-derivative}
 \partial_{s_{ab}}
      =\sum_{i\in S_a,\;j\in S_b}\partial_{z_{ij}},
\end{equation}
where $z_{ij}$ multiplies just $C_{ij}$ and $C_{ji}$.
Choose one endpoint pair for each tree edge and let $d_i$ count
the resulting degree of cell $i$. No chosen endpoint pair repeats.
At the real forest mask, allow complex increments with radii
\begin{equation}\label{oi-radii}
 r_{ij}=\frac{\delta}{\max(d_i,d_j)J_{ij}}.
\end{equation}
Zero covariance edges contribute zero and can be omitted.
Every block row and column of the perturbation has sum at most
$\delta$, hence its operator norm is at most $\delta$. On the
whole polydisc,
\[
 \operatorname{Herm}M\ge\lambda I/2,\quad
 \|M\|\le2\Lambda_0,\quad
 \operatorname{Herm}M^{-1}\ge V_*^{-1}I,\quad
 |\det M|\ge(\lambda/2)^{D|\cup S_a|}.
\]
Thus $|p_M|$ is bounded by $K_G^{|\cup S_a|}$ times the product
real Gaussian density with cell covariance $V_*I$. For singletons
use $|f_i|\le M_i+\mathbf1_{\mathrm{bad},i}$ and its independent
Gaussian tail bound $p_R$. For other supports use $|f_S|\le M_S$.
This proves
\[
 \left|\mathbb E_M\prod_a f_{S_a}\right|
       \le\prod_a b(S_a).
\]
Cauchy's estimate, once in every chosen covariance edge, yields
\begin{equation}\label{oi-tree-bound}
 \left|\partial_{T,\ {\rm endpoints}}
          \mathbb E_{C(s)}\prod_a f_{S_a}\right|
 \le c_G^{n-1}\prod_a b(S_a)
                   \prod_i d_i!\prod_{ij\in T}J_{ij}.
\end{equation}
Indeed,
$\prod_{ij}\max(d_i,d_j)\le\prod_i d_i^{d_i}
 \le e^{2(n-1)}\prod_i d_i!$.
The factorials are attached to original cells, including degree-zero
cells with $0!=1$, and not to the sizes or degrees of whole molecules.
Neither indicators nor interaction multipliers are differentiated.

Forest components now regroup \eqref{oi-disjoint} into the exact
hard-core gas
\begin{equation}\label{oi-effective-gas}
 Z_V=\sum_{\mathcal Q\ {\rm pairwise\ disjoint}}
                                 \prod_{X\in\mathcal Q}\zeta(X).
\end{equation}
Here $\zeta(X)$ sums partitions of $X$ into disjoint primitive
supports and molecule trees with their BKAR integrals. The
one-primitive case is its un-interpolated expectation.

\subsection{Endpoint factorials and the rooted-tree sum}

\begin{lemma}\label{oi-tree-lemma}
Under \eqref{oi-primitive-norm} and $\eta\le\eta_*$,
\begin{equation}\label{oi-zeta-bound}
 \sup_i\sum_{X\ni i}|\zeta(X)|
                  e^{\kappa\ell(X)+2|X|}\le\eta .
\end{equation}
\end{lemma}
\begin{proof}
Put $\widehat b(S)=b(S)e^{\kappa\ell(S)+2|S|}$ and
$\widehat J_{ij}=e^{\kappa d(i,j)}J_{ij}$.
The union weight is dominated by these support weights and endpoint
edge weights: spanning trees inside the supports together with
the chosen covariance edges connect their union.

Root a molecule tree and retain the incoming cell at each nonroot
vertex. Let $F_i$ majorize planted trees whose root molecule
contains $i$ and carries one extra incoming half-edge at $i$.
This half-edge has no covariance weight; it increments the cell
degree. Allowing repeated or overlapping primitive supports
enlarges the positive majorant. With
$t_k=c_G\sum_j\widehat J_{kj}F_j$, its recursion is
\begin{equation}\label{oi-recursion}
 F_i=\sum_{S\ni i}\widehat b(S)
              (1-t_i)^{-1}\prod_{k\in S}(1-t_k)^{-1}.
\end{equation}
To verify its normalization, suppose $n_k$ children depart from
cell $k$. The unordered-child factor, after summing endpoint
assignments, is $1/\prod_k n_k!$. The cell-degree factorial in
\eqref{oi-tree-bound} is $\prod_k n_k!$, except that the incoming
cell has $(n_i+1)!$. The residual factors are $1$ at every
ordinary cell and $n_i+1$ at the incoming cell. Their nonnegative
geometric series give exactly \eqref{oi-recursion}.
This can equivalently be obtained from ordered children weighted
by the inverse factorial of their total number. It is the rooted
labelled-tree exponential series; repeated labels are used only
for the upper bound.

Define the recursion by finite-height truncation. If $F_j\le2\eta$,
then $t_k\le2c_GA_\kappa\eta\le1/2$, and
\[
 F_i\le2\sum_{S\ni i}\widehat b(S)2^{|S|}
      \le2\sum_{S\ni i}\widehat b(S)e^{|S|}
      \le2\eta.
\]
This gives a convergent planted majorant. At an unplanted root
the extra factor $(1-t_i)^{-1}$ is absent, so its sum over root
supports containing $i$ is at most
$\sum_{S\ni i}\widehat b(S)2^{|S|}\le\eta$.
In an original disjoint primitive family whose union contains
$i$, precisely one support contains $i$; root the tree there.
This proves \eqref{oi-zeta-bound}.

In particular a molecule of $m$ cells with equal endpoint parameter
$t$ has unplanted generating function $(1-t)^{-m}$. Its degree-$d$
coefficient is $(m)_d/d!$, a rising factorial divided by $d!$.
Replacing this by the crude block bound $m^d d!$ would lose the
support-exponential control.
\end{proof}

\subsection{Logarithm and analytic comparison}

For $z'(X)=|\zeta(X)|e^{\kappa\ell(X)+|X|}$,
\eqref{oi-zeta-bound} gives the Kotecky--Preiss condition
with size function $|X|$:
\[
 \sum_{Y:\,Y\cap X\ne\varnothing}z'(Y)e^{|Y|}
       \le |X|\eta\le |X|.
\]
The absolutely convergent Mayer logarithm, grouped by the union
of each connected overlap cluster, defines $U_V$. The union
weight is at most the product of its polymer weights. Anchor a
cluster at $i$ by distinguishing a polymer containing $i$ and
apply the pinned cluster estimate. It gives
\begin{equation}\label{oi-pinned}
 \sup_i\sum_{X\ni i}e^{\kappa\ell(X)+|X|}|U_V(X)|
 \le\sup_i\sum_{Y\ni i}
       |\zeta(Y)|e^{\kappa\ell(Y)+2|Y|}\le\eta.
\end{equation}
Uniform background suprema can be taken at every preceding step.
Locally uniform convergence proves holomorphy and identifies
$Z_V$ with the exponential of this logarithm, proving
nonvanishing and \eqref{oi-log}.

For \eqref{oi-difference}, replace every $w_Y$ by $t w_Y$ while
retaining the bad singleton marks. This interpolates bounded
multipliers. Its interaction norm is $|t|\epsilon$.
If $\epsilon>0$, set $\rho=\eta_*/(2\epsilon)\ge2$.
Under \eqref{oi-quarter}, all estimates hold for $|t|\le\rho$
because $\eta(t)\le3\eta_*/4$. The canonical activity vector is
holomorphic on this disk in the rooted activity norm and bounded
by $2\eta_*$. Banach-valued Cauchy therefore gives
\[
 \|U(1)-U(0)\|\le\frac{2\eta_*}{\rho-1}\le8\epsilon .
\]
The case $\epsilon=0$ is immediate. At $t=0$ only bad singleton
primitives remain, so $U(0)$ is the original pure Gaussian cutoff
forest/Mayer expansion. This completes the proof.

\subsection{Scope of the abstract result}

The same conclusion holds for countably many labelled interactions
under \eqref{oi-smallness}, defining the integrand by the uniformly
convergent product of their bounded multipliers on the cutoff,
or equivalently by its finite-family limit. No convergence of the
chosen logarithms $V_Y$ themselves is asserted. If
$u_Y=\sup_{\rm cutoff}|V_Y|\le1$, then
$v_Y\le2u_Y$ is a sufficient conversion to the required norm.
Finite range is optional; the weighted summability is the actual
hypothesis. The factor $K_G^{|Y|}$ budgets the Gaussian density
cost but does not allow the original cell dimension $D$ to grow
with the regulator.

The argument establishes locality of support labels in fluctuation
space. A covariance depending on a remote background variable may
give every activity the same remote dependence, even if the
covariance is diagonal. Strict background support consequently
requires the separate local-precision/support-algebra hypotheses.
Actual Yang--Mills use also requires the stated independent cell
coordinates, uniform complex spectrum, fixed cutoff domain and
weighted interaction bounds. None is inferred here from a global
interaction stability estimate or from pointwise inverse bounds.

The imported expansion results are the BKAR forest formula on the
convex hull of partition masks (Abdesselam--Rivasseau,
arXiv:hep-th/9409094) and the Kotecky--Preiss hard-core polymer
theorem, \emph{Communications in Mathematical Physics} \textbf{103}
(1986), 491--498. The molecule grouping, cell-endpoint Cauchy
estimate and endpoint-factorial tree sum above supply the new
overlapping-support step.

\end{document}

